\section{Conclusion}\label{a:sec:conclusion}

The constant \(A(p)\) is well defined, bounded and computable, and its reciprocal
is the quasi-stationary mean of subcritical binary Galton--Watson branching.  The
product and series representations expose complementary features of the survival
recursion, while the near-critical theorem fixes its endpoint scale as
\(A(p)=2\varepsilon[1+O(\varepsilon\log(1/\varepsilon))]\).

The dynamical identification
\(A(p)=2\psi_{2p}(\tfrac12)\), interpreted as a basin value, explains why the
analytic problem is governed by Schr\"oder linearisation.  For every fixed
\(r\in(0,1)\), the function \(z\mapsto\psi_r(z)\) is
hypertranscendental over \(\mathbb C(z)\); the attracting window has no exceptional
parameter because the differentially algebraic cases classified by
Becker--Bergweiler occur only over repelling fixed points.

This theorem has a precise boundary.  It proves neither irrationality nor
transcendence of an individual \(A(p_0)\), and it does not show that
\(p\mapsto A(p)\) is non-elementary or differentially transcendental in \(p\).
The eleven-parameter PSLQ study supplies only finite-height negative evidence.
Thus the analytic obstruction is established for the conjugacy in its dynamical
variable, while the corresponding value and parameter questions remain open.
